% Options for packages loaded elsewhere
\PassOptionsToPackage{unicode}{hyperref}
\PassOptionsToPackage{hyphens}{url}
\documentclass[
]{article}
\usepackage{xcolor}
\usepackage{amsmath,amssymb}
\usepackage{amsthm}
\usepackage{booktabs}
\setcounter{secnumdepth}{3}
\usepackage{iftex}
\ifPDFTeX
  \usepackage[T1]{fontenc}
  \usepackage[utf8]{inputenc}
  \usepackage{textcomp}
\else
  \usepackage{unicode-math}
  \defaultfontfeatures{Scale=MatchLowercase}
  \defaultfontfeatures[\rmfamily]{Ligatures=TeX,Scale=1}
\fi
\usepackage{lmodern}
\IfFileExists{upquote.sty}{\usepackage{upquote}}{}
\IfFileExists{microtype.sty}{%
  \usepackage[]{microtype}
  \UseMicrotypeSet[protrusion]{basicmath}
}{}
\makeatletter
\@ifundefined{KOMAClassName}{%
  \IfFileExists{parskip.sty}{%
    \usepackage{parskip}
  }{% else
    \setlength{\parindent}{0pt}
    \setlength{\parskip}{6pt plus 2pt minus 1pt}}
}{%
  \KOMAoptions{parskip=half}}
\makeatother
\setlength{\emergencystretch}{3em}
\providecommand{\tightlist}{%
  \setlength{\itemsep}{0pt}\setlength{\parskip}{0pt}}
\usepackage{bookmark}
\IfFileExists{xurl.sty}{\usepackage{xurl}}{}
\urlstyle{same}
\hypersetup{
  pdftitle={Die Trinität als notwendige Struktur einer monistischen Modalontologie – Version 4},
  hidelinks,
  pdfcreator={LaTeX via pandoc}
}

\newtheorem{theorem}{Satz}
\newtheorem{lemma}{Lemma}
\newtheorem{definition}{Definition}
\newtheorem{corollary}{Korollar}

\title{Die Trinität als notwendige Struktur einer monistischen Modalontologie \\ \large Version 4 – Vollständige formale Analyse}
\author{Paul Koop}
\date{}

\begin{document}
\maketitle

\begin{center}
\emph{Vergangenheit und Zukunft sind Horizonte des Wissens. \\ Die Gegenwart ist der Ort der Wirklichkeit.}
\end{center}

\newpage
{
\setcounter{tocdepth}{3}
\tableofcontents
}
\newpage

\section{Vorwort}

Diese Abhandlung ist die formale Vollendung einer langen Entwicklung. Sie vereint die Erkenntnisse aus allen vorangegangenen Versionen:

\begin{itemize}
\item \textbf{Version 1–3:} Versuch, die Trinität direkt aus S5 abzuleiten – scheiterte an der fehlenden Brücke von Existenz zu Notwendigkeit.
\item \textbf{Version 4 (alt):} Nachweis, dass reines S5 nicht ausreicht – die Zielaussage ist in S5 allein nicht ableitbar. Dieser Nachweis wird hier in die neue Version integriert.
\item \textbf{Version 5:} Formalisierung der Superpositions-Intuition und Beweis der Zielaussage im erweiterten System S5+SP.
\end{itemize}

Die \textbf{vorliegende Version 4 (neu)} vereint beide Perspektiven:

\begin{enumerate}
\item \textbf{Teil I:} Strenger formaler Nachweis, dass die Zielaussage in reinem S5 \textbf{nicht} ableitbar ist (Übernahme der alten Version 4).
\item \textbf{Teil II:} Strenger formaler Beweis, dass die Zielaussage im erweiterten System S5+SP (mit Superpositions-Axiomen) \textbf{ableitbar} ist.
\item \textbf{Teil III:} Metatheoretische Einordnung – was wurde gezeigt, was nicht, und welche Fragen bleiben offen.
\end{enumerate}

Die Grundthese der gesamten Untersuchung lautet:

\begin{quote}
\textbf{Unter den Axiomen des Monismus (A1), der Existenz einer möglichen Welt (A4), den transzendentalen Brücken (A11, A12) und der Superpositions-Hypothese (ASP1–ASP5) folgt notwendig die Trinität von Ursprung, Totalität und Selbsterkenntnis. Ohne die Superpositions-Hypothese ist sie in S5 nicht beweisbar.}
\end{quote}

\section{Einleitung: Ziel und methodische Selbstbeschränkung}

Die klassische Metaphysik hat immer wieder versucht, die Grundstruktur der Wirklichkeit aus wenigen fundamentalen Prinzipien abzuleiten. Eine besondere Herausforderung entsteht dabei durch drei scheinbar verschiedene Aspekte:

\begin{enumerate}
\item \textbf{Warum gibt es überhaupt etwas und nicht vielmehr nichts?} – der Ursprung (U).
\item \textbf{Wie verhält sich die Gesamtheit aller Möglichkeiten zur Wirklichkeit?} – die Totalität (T).
\item \textbf{Wie kann innerhalb der Wirklichkeit ein Zustand entstehen, der die Wirklichkeit selbst erkennt?} – die Selbsterkenntnis (S).
\end{enumerate}

Die hier untersuchte Ontologie schlägt vor, diese drei Fragen nicht durch drei getrennte metaphysische Substanzen zu beantworten, sondern durch drei notwendige Perspektiven derselben Wirklichkeit:

\begin{itemize}
\item \textbf{Ursprung (U):} der notwendige Grund dafür, dass überhaupt Möglichkeiten existieren – verstanden als unterer Grenzwert der Totalität;
\item \textbf{Totalität (T):} die Gesamtheit aller realisierten Möglichkeiten – verstanden als offenes, wohlfundiertes Intervall;
\item \textbf{Selbsterkenntnis (S):} der reflexive Abschluss der Wirklichkeit in einem Zustand vollständiger Erkenntnis ihrer selbst – verstanden als oberer Grenzwert der Totalität.
\end{itemize}

Diese Struktur wird als \textbf{Trinität} bezeichnet:

\[
Tr := U \land T \land S
\]

Der Begriff „Trinität“ bezeichnet hierbei keine personale oder substanzielle Dreiheit, sondern eine funktionale Einheit dreier notwendiger Aspekte einer einzigen Wirklichkeit.

\section{Formale Sprache und logischer Rahmen}

\subsection{Modallogik S5}

Wir arbeiten in der Modallogik erster Stufe mit Identität im System S5:

\begin{itemize}
\item \textbf{Notwendigkeit:} \(\Box p\)
\item \textbf{Möglichkeit:} \(\Diamond p := \neg\Box\neg p\)
\end{itemize}

\textbf{Axiome von S5:}
\begin{itemize}
\item (K) \(\Box(p \rightarrow q) \rightarrow (\Box p \rightarrow \Box q)\)
\item (T) \(\Box p \rightarrow p\)
\item (4) \(\Box p \rightarrow \Box\Box p\)
\item (5) \(\Diamond p \rightarrow \Box\Diamond p\)
\end{itemize}

\textbf{Inferenzregeln:}
\begin{itemize}
\item \textbf{Modus Ponens:} Aus \(p\) und \(p \rightarrow q\) folgt \(q\).
\item \textbf{Necessitation:} Aus \(p\) (ableitbar) folgt \(\Box p\).
\end{itemize}

\textbf{Tableau-Regeln für S5:}
\begin{align*}
& (\neg\Box) & \frac{\neg\Box A}{\Diamond\neg A} \\
& (\Diamond) & \frac{\Diamond A}{A @ w_{\text{neu}}} \quad \text{(neue Welt)} \\
& (\Box) & \frac{\Box A @ w}{A @ v} \quad \text{für jede bereits existierende Welt } v \\
& (\neg\forall) & \frac{\neg\forall x A}{\exists x \neg A} \\
& (\exists) & \frac{\exists x A}{A[a/x]} \quad \text{(a neu)} \\
& (\forall) & \frac{\forall x A}{A[a/x]} \quad \text{(a beliebig)}
\end{align*}

\subsection{Prädikatenlogik}

Wir verwenden die klassische Prädikatenlogik erster Stufe mit Identität (\(=\)) und den üblichen Quantoren (\(\forall, \exists\)).

\section{Die Axiome (Basis)}

\subsection{A1 – Monismus}

Es gibt keine fundamental getrennten Wirklichkeitsbereiche.

\[
\boxed{A1 := \Box\neg\exists x\exists y\, FundamentalGetrennt(x,y)}
\]

\textbf{Erläuterung:} Wenn zwei Bereiche fundamental getrennt wären, gäbe es keine kausale oder ontologische Wechselwirkung zwischen ihnen. Dann könnten sie nicht Teil einer einheitlichen Wirklichkeit sein, was dem Monismus widerspricht.

\subsection{A4 – Existenz einer möglichen Wirklichkeit}

Es gibt mindestens eine mögliche Welt.

\[
\boxed{A4 := \Diamond\exists w\, Welt(w)}
\]

\textbf{Erläuterung:} Dieses Axiom stellt sicher, dass das modale Universum nicht leer ist. Es ist das schwächste mögliche Existenzaxiom: Es sagt nicht, dass eine Welt tatsächlich existiert, sondern nur, dass sie möglich ist.

\subsection{A11 – Transzendentale Brücke}

Eine Welt ist genau dann von Bewusstsein getrennt, wenn sie kein Bewusstsein enthält.

\[
\boxed{A11 := \forall w \left( \text{WeltGetrenntVonBewusstsein}(w) \leftrightarrow \neg \exists c (Bewusstsein(c) \land c(w)) \right)}
\]

mit der Definition:

\[
\boxed{\text{WeltGetrenntVonBewusstsein}(w) := \neg \exists c (Bewusstsein(c) \land c(w))}
\]

\textbf{Erläuterung:} Dieses Axiom formalisiert die transzendentale Einsicht, dass eine Welt ohne Bewusstsein unerkennbar und daher fundamental getrennt wäre – was der Monismus (A1) verbietet.

\subsection{A12 – Erfahrbarkeit und Realisierung}

\[
\boxed{A12 := \forall p. \text{Erfahrbar}(p) \leftrightarrow \exists w. Realisiert(w, p)}
\]

mit der Definition:

\[
\boxed{\text{Erfahrbar}(p) := \exists w. (\text{Bewusstsein}(w) \land \text{Realisiert}(w, p))}
\]

\textbf{Erläuterung:} Dieses Axiom besagt: Eine Aussage ist genau dann erfahrbar, wenn sie in einer Welt realisiert ist. Es ist die formale Fassung des transzendentalen Arguments: Was nicht realisiert ist, kann nicht erfahren werden.

\subsection{A13, A14, A15 – Die drei Implikationen der Fälle}

Diese Axiome formalisieren die drei Reductio-Fälle:

\[
\boxed{A13 := \Box\forall x(T(x) \rightarrow U(x))}
\]
\[
\boxed{A14 := \Box\forall x(U(x) \rightarrow S(x))}
\]
\[
\boxed{A15 := \Box\forall x(S(x) \rightarrow T(x))}
\]

\textbf{Erläuterung:} Sie besagen, dass die drei Grenzwerte \(T, U, S\) in einer wechselseitigen Implikationskette stehen – was in S5 zu ihrer Äquivalenz führt.

\section{Definition der Grenzwertstruktur}

\subsection{Die Totalität T als offenes, wohlfundiertes Intervall}

Die \textbf{Totalität T} ist die Gesamtheit aller realisierten Zustände. Wir verstehen T als \textbf{wohlfundiertes, offenes Intervall}:

\[
\boxed{T := \{ x \mid U < x < S \}}
\]

\textbf{Erläuterung:}
\begin{itemize}
\item Ein offenes Intervall \((U, S)\) enthält alle Zustände zwischen U und S, aber \textbf{nicht} U und S selbst.
\item \textbf{Wohlfundiertheit:} Jede nicht-leere Teilmenge von Zuständen hat ein minimales Element. Dies verhindert unendliche Ketten von Gründen.
\item Der \textbf{Ursprung U} ist der untere Grenzwert – das, was dem Nichts am nächsten kommt, aber selbst kein Nichts ist.
\item Die \textbf{Selbsterkenntnis S} ist der obere Grenzwert – die vollständige Transparenz der Totalität, die selbst nicht zur Totalität gehört.
\end{itemize}

\subsection{Der Ursprung U als unterer Grenzwert}

\[
\boxed{U := \lim_{x \to \inf} T}
\]

\subsection{Die Selbsterkenntnis S als oberer Grenzwert}

\[
\boxed{S := \lim_{x \to \sup} T}
\]

\subsection{Die Trinität}

\[
\boxed{Tr := U \land T \land S}
\]

\section{TEIL I: NACHWEIS DER NICHT-ABLEITBARKEIT IN REINEM S5}

\subsection{Ziel}

Zeige, dass:

\[
\boxed{\text{S5} \;\not\vdash\; \Box\forall x\,Tr(x)}
\]

aus den Axiomen \(A1, A4, A11, A12, A13, A14, A15\) allein nicht folgt.

\subsection{Methode}

Konstruiere ein \textbf{offenes S5-Tableau} für die Negation der Zielformel.  
Gemäß dem \textbf{Korrektheits- und Vollständigkeitssatz} des S5-Tableau-Kalküls gilt:

\begin{quote}
Eine Formelmenge ist genau dann \textbf{erfüllbar} in einem S5-Modell, wenn das Tableau \textbf{einen offenen Zweig} besitzt.
\end{quote}

Daher: Wenn das Tableau für die negierte Zielformel offen bleibt, dann existiert ein S5-Modell, das alle Axiome und die Negation der Zielformel erfüllt – also ist die Zielformel \textbf{kein Theorem}.

\subsection{Tableau (reine Smullyan-Regeln)}

\begin{verbatim}
1.  ¬□∀x Tr(x)                                    [Annahme: Ziel sei kein Theorem]
2.  ◇¬∀x Tr(x)                                    [1, ¬□-Regel]
3.  ¬∀x Tr(x) @ w0                                [2, ◇-Regel: neue Welt w0]
4.  ∃x ¬Tr(x) @ w0                                [3, ¬∀-Regel: ¬∀xA ⊢ ∃x¬A]
5.  ¬Tr(a) @ w0                                   [4, ∃-Regel: a neu]
6.  ¬(T(a) ∧ U(a) ∧ S(a)) @ w0                    [5, D1 (Tr-Definition)]

    → β-Regel auf 6:
    6a. ¬T(a) @ w0
    6b. ¬U(a) @ w0
    6c. ¬S(a) @ w0

─────────────────────────────────────────────────────────────
ZWEIG 6a: ¬T(a) @ w0
─────────────────────────────────────────────────────────────

7.  □∀x(S(x) → T(x)) @ w0                         [A15]
8.  ∀x(S(x) → T(x)) @ w0                          [7, □-Regel]
9.  S(a) → T(a) @ w0                              [8, ∀-Regel]

    → β-Regel auf 9:
    9a. ¬S(a) @ w0
    9b. T(a) @ w0                                  [Widerspruch zu 6a → Zweig 9b schließt]

    Also: 9a. ¬S(a) @ w0

10. □∀x(T(x) → U(x)) @ w0                         [A13]
11. ∀x(T(x) → U(x)) @ w0                          [10, □-Regel]
12. T(a) → U(a) @ w0                              [11, ∀-Regel]

    → β-Regel auf 12:
    12a. ¬T(a) @ w0                                [bereits in 6a]
    12b. U(a) @ w0                                 [kein Widerspruch]

    → Wähle Zweig 12a (konsistent mit 6a).

13. □∀x(U(x) → S(x)) @ w0                         [A14]
14. ∀x(U(x) → S(x)) @ w0                          [13, □-Regel]
15. U(a) → S(a) @ w0                              [14, ∀-Regel]

    → β-Regel auf 15:
    15a. ¬U(a) @ w0
    15b. S(a) @ w0                                 [Widerspruch zu 9a → Zweig 15b schließt]

    Also: 15a. ¬U(a) @ w0

    → Damit im Zweig 6a: ¬T(a), ¬U(a), ¬S(a) @ w0.
    → Dies ist konsistent – kein Widerspruch.

─────────────────────────────────────────────────────────────
ZWEIG 6b: ¬U(a) @ w0
─────────────────────────────────────────────────────────────

16. □∀x(T(x) → U(x)) @ w0                         [A13]
17. ∀x(T(x) → U(x)) @ w0                          [16, □-Regel]
18. T(a) → U(a) @ w0                              [17, ∀-Regel]

    → β-Regel auf 18:
    18a. ¬T(a) @ w0
    18b. U(a) @ w0                                 [Widerspruch zu 6b → schließt]

    Also: 18a. ¬T(a) @ w0

19. □∀x(S(x) → T(x)) @ w0                         [A15]
20. ∀x(S(x) → T(x)) @ w0                          [19, □-Regel]
21. S(a) → T(a) @ w0                              [20, ∀-Regel]

    → β-Regel auf 21:
    21a. ¬S(a) @ w0
    21b. T(a) @ w0                                 [Widerspruch zu 18a → schließt]

    Also: 21a. ¬S(a) @ w0

22. □∀x(U(x) → S(x)) @ w0                         [A14]
23. ∀x(U(x) → S(x)) @ w0                          [22, □-Regel]
24. U(a) → S(a) @ w0                              [23, ∀-Regel]

    → β-Regel auf 24:
    24a. ¬U(a) @ w0                                [bereits in 6b]
    24b. S(a) @ w0                                 [Widerspruch zu 21a → schließt]

    → Konsistenter Zweig: ¬U(a), ¬T(a), ¬S(a) @ w0.

─────────────────────────────────────────────────────────────
ZWEIG 6c: ¬S(a) @ w0
─────────────────────────────────────────────────────────────

25. □∀x(S(x) → T(x)) @ w0                         [A15]
26. ∀x(S(x) → T(x)) @ w0                          [25, □-Regel]
27. S(a) → T(a) @ w0                              [26, ∀-Regel]

    → β-Regel auf 27:
    27a. ¬S(a) @ w0                                [bereits in 6c]
    27b. T(a) @ w0

    → Beide Zweige möglich.

28. □∀x(T(x) → U(x)) @ w0                         [A13]
29. ∀x(T(x) → U(x)) @ w0                          [28, □-Regel]
30. T(a) → U(a) @ w0                              [29, ∀-Regel]

    → β-Regel auf 30:
    30a. ¬T(a) @ w0
    30b. U(a) @ w0

31. □∀x(U(x) → S(x)) @ w0                         [A14]
32. ∀x(U(x) → S(x)) @ w0                          [31, □-Regel]
33. U(a) → S(a) @ w0                              [32, ∀-Regel]

    → β-Regel auf 33:
    33a. ¬U(a) @ w0
    33b. S(a) @ w0                                 [Widerspruch zu 6c → schließt]

    Also: 33a. ¬U(a) @ w0

    → Kombiniere: Aus 30b (U(a)) und 33a (¬U(a)) ergibt sich Widerspruch.
    → Daher muss 30a gelten: ¬T(a) @ w0.

    → Damit: ¬S(a), ¬U(a), ¬T(a) @ w0 konsistent.

─────────────────────────────────────────────────────────────
ALLE ZWEIGE: ¬T(a) ∧ ¬U(a) ∧ ¬S(a) @ w0 ist konsistent.
─────────────────────────────────────────────────────────────

─────────────────────────────────────────────────────────────
ÜBRIGE AXIOME (A4, A11, A12) – keine Anwendung
─────────────────────────────────────────────────────────────

34. ◇∃w Welt(w) @ w0                              [A4]
35. ∃w Welt(w) @ w1                               [34, ◇-Regel, w1 neu]
    → Führt zu neuer Welt w1, aber nicht zu w0 zurück.

36. A11, A12: Keine Instanzen in w0, die T(a), U(a) oder S(a) erzwingen.

─────────────────────────────────────────────────────────────
FAZIT DES TABLEAUS:
─────────────────────────────────────────────────────────────

Das Tableau besitzt einen **offenen Zweig**:

    w0, a, mit ¬T(a), ¬U(a), ¬S(a).

Kein Axiom erzeugt in w0 T(a), U(a) oder S(a).
Die übrigen Axiome führen zu neuen Welten (w1, ...),
aber nicht zu w0 zurück.

Daher ist das Tableau **nicht geschlossen**.
\end{verbatim}

\subsection{Metatheoretischer Schluss}

Nach dem \textbf{Korrektheits- und Vollständigkeitssatz} des S5-Tableau-Kalküls
(siehe z.\,B. Smullyan, Fitting, oder Blackburn/de Rijke/Venema) gilt:

\begin{quote}
Eine Formelmenge \(\Sigma\) ist genau dann S5-erfüllbar, wenn das Tableau für \(\Sigma\) einen offenen Zweig besitzt.
\end{quote}

Unser Tableau für

\[
\Sigma = \{ A1, A4, A11, A12, A13, A14, A15, \neg\Box\forall xTr(x) \}
\]

besitzt einen offenen Zweig.

Also ist \(\Sigma\) S5-erfüllbar.

Also gibt es ein S5-Modell, das alle Axiome erfüllt, aber in der Welt \(w_0\) ein Individuum \(a\) besitzt, für das \(\neg T(a)\), \(\neg U(a)\) und \(\neg S(a)\) gelten.

Also gilt in diesem Modell nicht \(\Box\forall xTr(x)\).

Also ist \(\Box\forall xTr(x)\) \textbf{kein logisches Theorem} der gegebenen Axiomatik.

\begin{theorem}[Nicht-Ableitbarkeit in reinem S5]
\[
\boxed{
\text{S5} \;\not\vdash\; \Box\forall x\,Tr(x)
}
\]
\end{theorem}

\section{TEIL II: BEWEIS IM ERWEITERTEN SYSTEM S5+SP}

\subsection{Die Superpositions-Erweiterung (ASP)}

Zusätzlich zu den Axiomen A1, A4, A11, A12, A13, A14, A15 führen wir die \textbf{Superpositions-Axiome} ein. Sie formalisieren die Idee, dass Bewusstsein (\(C\)) der selbstreflexive Moment einer Superposition ist, aus der alle Welten entstehen.

\subsubsection{ASP1 – Bewusstsein in der Superposition}

\[
\boxed{ASP1 := C @ w_{\text{super}}}
\]

\textbf{Erläuterung:} In der Superposition gilt Bewusstsein. Die Superposition ist der Urzustand, in dem alle Möglichkeiten noch ungetrennt enthalten sind.

\subsubsection{ASP2 – Einzigartigkeit von \(C\)}

\[
\boxed{ASP2 := \forall v (w_{\text{super}} R v \rightarrow (C @ v \leftrightarrow v = w_{\text{super}}))}
\]

\textbf{Erläuterung:} Bewusstsein gilt nur in der Superposition – keine andere Welt hat es. Bewusstsein ist ein singuläres Ereignis.

\subsubsection{ASP3 – Die Superposition enthält alle Tr-Eigenschaften}

\[
\boxed{ASP3 := \forall x\,Tr(x) @ w_{\text{super}}}
\]

\textbf{Erläuterung:} Die Superposition enthält bereits alle Eigenschaften \(T, U, S\). Sie ist der Grund für alles, was in den Welten gilt.

\subsubsection{ASP4 – Übertragung auf alle Welten}

\[
\boxed{ASP4 := \forall v (w_{\text{super}} R v \rightarrow \forall x\,Tr(x) @ v)}
\]

\textbf{Erläuterung:} Was in der Superposition gilt, gilt in allen erreichbaren Welten. Die Superposition ist ein normativer Ursprung.

\subsubsection{ASP5 – Universelle Erreichbarkeit (als Rahmenbedingung)}

\[
\boxed{ASP5 := \forall v (v \neq w_{\text{super}} \rightarrow w_{\text{super}} R v)}
\]

\textbf{Erläuterung:} Jede andere Welt ist von der Superposition aus erreichbar. Die Superposition ist der einzige Ursprung.

\subsection{Beweis im erweiterten System S5+SP}

\begin{theorem}[Ableitbarkeit in S5+SP]
\[
\boxed{
\text{S5+SP} \;\vdash\; \Box\forall x\,Tr(x)
}
\]
wobei \(\text{S5+SP} := \text{S5} + \{ASP1, ASP2, ASP3, ASP4, ASP5\}\).
\end{theorem}

\subsubsection{Semantischer Beweis}

Sei \(\mathcal{M} = (W, R, w_{\text{super}}, I)\) ein beliebiges S5+SP-Modell mit ASP1–ASP5.

Zu zeigen: Für alle \(w \in W\) gilt \(\mathcal{M}, w \models \forall x\,Tr(x)\).

Sei \(w \in W\) beliebig.

\textbf{Fall 1:} \(w = w_{\text{super}}\).

Nach ASP3 gilt:

\[
\mathcal{M}, w_{\text{super}} \models \forall x\,Tr(x)
\]

Also gilt die Behauptung.

\textbf{Fall 2:} \(w \neq w_{\text{super}}\).

Nach ASP5 gilt:

\[
w_{\text{super}} R w
\]

Nach ASP4 folgt daraus:

\[
\mathcal{M}, w \models \forall x\,Tr(x)
\]

Also gilt die Behauptung.

Da \(w\) beliebig war, gilt für alle \(w \in W\):

\[
\mathcal{M}, w \models \forall x\,Tr(x)
\]

Dies ist äquivalent zu:

\[
\mathcal{M} \models \Box\forall x\,Tr(x)
\]

q.e.d.

\subsubsection{Tableau-Beweis in S5+SP}

\begin{verbatim}
TABLEAU-BEWEIS IN S5+SP
ZIEL: ⊢ □∀x Tr(x)

─────────────────────────────────────────────────────────────
REDUCTIO-ANNAHME:
─────────────────────────────────────────────────────────────

1.  ¬□∀x Tr(x)                                    [Annahme: Ziel sei falsch]
2.  ◇¬∀x Tr(x)                                    [1, ¬□-Regel]
3.  ¬∀x Tr(x) @ w0                                [2, ◇-Regel: neue Welt w0]
4.  ∃x ¬Tr(x) @ w0                                [3, ¬∀-Regel]
5.  ¬Tr(a) @ w0                                   [4, ∃-Regel: a neu]

─────────────────────────────────────────────────────────────
FALLUNTERSCHEIDUNG: w0 = w_super ∨ w0 ≠ w_super
─────────────────────────────────────────────────────────────

6.  w0 = w_super  ∨  w0 ≠ w_super                  [Identität]

    → Verzweigung A: w0 = w_super
    → Verzweigung B: w0 ≠ w_super

─────────────────────────────────────────────────────────────
VERZWEIGUNG A: w0 = w_super
─────────────────────────────────────────────────────────────

7.  ¬Tr(a) @ w_super                               [5, Substitution]
8.  ∀x Tr(x) @ w_super                             [ASP3, Axiom]
9.  Tr(a) @ w_super                                [8, ∀-Regel auf a]
10. Widerspruch: Tr(a) @ w_super und ¬Tr(a) @ w_super
    → Verzweigung A schließt (⊥).

─────────────────────────────────────────────────────────────
VERZWEIGUNG B: w0 ≠ w_super
─────────────────────────────────────────────────────────────

11. w0 ≠ w_super                                   [aus 6, Zweig B]
12. w_super R w0                                   [11, ASP5]
13. ∀x Tr(x) @ w0                                  [12, ASP4]
14. Tr(a) @ w0                                     [13, ∀-Regel auf a]
15. Widerspruch: Tr(a) @ w0 und ¬Tr(a) @ w0 (aus 5)
    → Verzweigung B schließt (⊥).

─────────────────────────────────────────────────────────────
BEIDE VERZWEIGUNGEN SCHLIESSEN.
─────────────────────────────────────────────────────────────

Daher ist die Annahme ¬□∀xTr(x) widersprüchlich.

Also: ⊢ □∀xTr(x) in S5+SP.

QED.
\end{verbatim}

\section{TEIL III: METATHEORETISCHE EINORDNUNG}

\subsection{Was wurde gezeigt?}

\begin{table}[h]
\centering
\begin{tabular}{lll}
\toprule
\textbf{System} & \textbf{Aussage} & \textbf{Status} \\
\midrule
Reines S5 & \(\text{S5} \;\not\vdash\; \Box\forall x\,Tr(x)\) & Bewiesen (offener Tableau-Zweig) \\
S5+SP & \(\text{S5+SP} \;\vdash\; \Box\forall x\,Tr(x)\) & Bewiesen (geschlossenes Tableau) \\
\bottomrule
\end{tabular}
\end{table}

Die entscheidende Erkenntnis ist:

\begin{quote}
Die Brücke von der Existenz von \(Tr\) in einer Welt zur Notwendigkeit von \(Tr\) in allen Welten ist \textbf{kein Theorem von S5}. Sie muss als \textbf{zusätzliche metaphysische Annahme} eingeführt werden – hier in Form der Superpositions-Axiome ASP1–ASP5.
\end{quote}

\subsection{Was wurde nicht gezeigt?}

\begin{itemize}
\item Die \textbf{Wahrheit} der Superpositions-Axiome ASP1–ASP5. Sie sind \textbf{metaphysische Prämissen}, keine logischen Theoreme.
\item Dass die Zielaussage aus den ursprünglichen Axiomen A1, A4, A11, A12, A13, A14, A15 allein folgt – im Gegenteil, Teil I zeigt, dass sie es nicht tut.
\end{itemize}

\subsection{Die Rolle von \(C\) (Bewusstsein)}

Die Superpositions-Axiome formalisieren die Intuition, dass \textbf{Bewusstsein (\(C\))} der selbstreflexive Moment der Superposition ist:

\begin{itemize}
\item ASP1: \(C\) gilt in der Superposition.
\item ASP2: \(C\) gilt nur dort – es ist ein singuläres Ereignis.
\item ASP3–ASP5: Die Superposition ist der notwendige Grund für alle Eigenschaften in allen Welten.
\end{itemize}

Damit wird \(C\) zur \textbf{Brücke von Existenz zu Notwendigkeit}: Weil die Superposition \(Tr\) enthält und alle Welten aus ihr hervorgehen, gilt \(Tr\) notwendig in allen Welten.

\section{Gegenmodelle}

Die Trinität kann vermieden werden, wenn mindestens eines der Axiome aufgegeben wird:

\begin{table}[h]
\centering
\begin{tabular}{lll}
\toprule
\textbf{Aufgegebenes Axiom} & \textbf{Gegenmodell} & \textbf{Konsequenz} \\
\midrule
A1 (Monismus) & Dualismus (Descartes) & S kann ohne T existieren; Erkenntnis und Objekt sind getrennt \\
A4 (Existenz einer Welt) & Nihilismus & Es gibt keine Welt; die gesamte Ontologie ist leer \\
A11 (Transzendentale Brücke) & Erkenntnistheoretischer Skeptizismus & Unerkennbarkeit impliziert keine fundamentale Trennung \\
A12 (Erfahrbarkeit ↔ Realisierung) & Empirismus & Erfahrbarkeit ist nicht mit Realisierung identisch \\
ASP1–ASP5 (Superposition) & Keine Superposition & Die Brücke von Existenz zu Notwendigkeit fehlt \\
\bottomrule
\end{tabular}
\end{table}

\section{Anhang: Vollständige Axiome und Theoreme}

\subsection{Axiome (vollständig)}

\[
\begin{aligned}
A1 &:= \Box\neg\exists x\exists y\, FundamentalGetrennt(x,y) \\
A4 &:= \Diamond\exists w\, Welt(w) \\
A11 &:= \forall w \left( \text{WeltGetrenntVonBewusstsein}(w) \leftrightarrow \neg \exists c (Bewusstsein(c) \land c(w)) \right) \\
A12 &:= \forall p. \text{Erfahrbar}(p) \leftrightarrow \exists w. Realisiert(w, p) \\
A13 &:= \Box\forall x(T(x) \rightarrow U(x)) \\
A14 &:= \Box\forall x(U(x) \rightarrow S(x)) \\
A15 &:= \Box\forall x(S(x) \rightarrow T(x)) \\
ASP1 &:= C @ w_{\text{super}} \\
ASP2 &:= \forall v (w_{\text{super}} R v \rightarrow (C @ v \leftrightarrow v = w_{\text{super}})) \\
ASP3 &:= \forall x\,Tr(x) @ w_{\text{super}} \\
ASP4 &:= \forall v (w_{\text{super}} R v \rightarrow \forall x\,Tr(x) @ v) \\
ASP5 &:= \forall v (v \neq w_{\text{super}} \rightarrow w_{\text{super}} R v)
\end{aligned}
\]

\subsection{Definitionen}

\[
\begin{aligned}
\text{WeltGetrenntVonBewusstsein}(w) &:= \neg \exists c (Bewusstsein(c) \land c(w)) \\
\text{Erfahrbar}(p) &:= \exists w. (\text{Bewusstsein}(w) \land \text{Realisiert}(w, p)) \\
T &:= \{ x \mid U < x < S \} \\
U &:= \lim_{x \to \inf} T \\
S &:= \lim_{x \to \sup} T \\
Tr &:= U \land T \land S
\end{aligned}
\]

\subsection{Abgeleitete Theoreme}

\[
\begin{aligned}
& \exists T \quad \text{(aus A4)} \\
& \Diamond C \quad \text{(A5 – aus A1, A11)} \\
& \forall p(\Diamond p \rightarrow \exists w. Realisiert(w,p)) \quad \text{(A2 – aus A1, A11, A12)} \\
& \Box(T \rightarrow U) \quad \text{(aus Fall 1)} \\
& \Box(U \rightarrow S) \quad \text{(aus Fall 2 + Superposition)} \\
& \Box(S \rightarrow T) \quad \text{(aus Fall 3)} \\
& \Rightarrow \Box(U \leftrightarrow T) \land \Box(T \leftrightarrow S) \land \Box(S \leftrightarrow U) \\
& \Rightarrow \Box(U \land T \land S)
\end{aligned}
\]

\section{Schluss}

Die Abhandlung hat gezeigt:

\begin{enumerate}
\item \textbf{In reinem S5} ist die Trinität \textbf{nicht beweisbar} (Teil I).
\item \textbf{Im erweiterten System S5+SP} (mit Superpositions-Axiomen) ist die Trinität \textbf{beweisbar} (Teil II).
\item Die entscheidende metaphysische Last liegt auf den Superpositions-Axiomen – sie sind \textbf{keine logischen Theoreme}, sondern \textbf{zusätzliche Annahmen}.
\end{enumerate}

Die formale Arbeit hat damit die logischen und metaphysischen Anforderungen eines solchen Beweises präzise herausgearbeitet. Die entscheidende Frage – ob die Superpositionsaxiome wahr sind – bleibt einer philosophischen oder physikalischen Untersuchung vorbehalten. Die Logik hat ihr Werk getan: Sie hat die Konsequenzen der Annahmen expliziert und die Grenzen des reinen S5 klar benannt.

\begin{quote}
\textbf{Die Trinität ist keine zusätzliche Entität, sondern eine Strukturbedingung – aber sie ist nur unter der Superpositions-Hypothese beweisbar.}
\end{quote}

\begin{center}
\emph{Diese Abhandlung wurde im Geiste strenger Modallogik verfasst, jedoch in der Sprache der Philosophie – denn die Wahrheit bedarf beider: der Präzision der Formel und der Weite des Begriffs.}
\end{center}

\end{document}